\chapter{Theory} \label{sec:Theory}
It is assumed that the reader has a fundamental understanding of electrical engineering principles. The purpose of this chapter is therefore to present and explain the theoretical concepts that are specifically relevant to this work and may not be familiar to all readers.

\section{Strain} \label{sec:Strain}
Strain is a measure of the relative deformation of a material subjected to mechanical loading. It describes how much a body stretches or compresses compared to its original dimensions. For a uniaxial deformation, the engineering strain is defined as:

\begin{align} \label{eq:strain definition}
    \varepsilon=\frac{\Delta L}{L_{0}},
\end{align}

where $\Delta L$ is the change in length and $L_0$ is the original length of the material. For small deformations, this definition provides an accurate and widely used description of material behavior~\cite{metal_gauge}.

During assembly of the bolted joint, the bolt undergoes tensile strain while the nut undergoes compressive strain, as shown in \Cref{fig:strain}. This deformation is directly related to the preload in the bolted joint, and a reduction in preload will result in a corresponding change in strain. This relationship enables strain to be used as an indicator for monitoring the mechanical state of the joint~\cite{Effects_of_preload_deficiency, StrainBand}.

\begin{figure}[h]
\centering
\includestandalone[mode=buildmissing]{Standalone/Drawings/strain-tikz}
    \caption{Illustration of strain in a bolted joint. The bolt undergoes tensile stretching ($\Delta L_b$), while the two nuts experience a corresponding compressive deformation ($\Delta L_n$), due to the applied preload.}
    \label{fig:strain}
\end{figure}